%% kapa: exam=1 round=20 year=2009 period=1 form=A subject=civil_law subject_ko=민법 seq=21 number=21
%% source: Q-Net 감정평가사 기출문제 (cst00309) artlSeq=5024178 file=extracted/1008/5024178/2009년 제20회 감정평가사 1차시험 기출문제/1교시.hwp
%% conversion: kordoc 4.24.51 -> tools/gen1.py (auto); review=proofread-sonnet; flags=none
\begin{question}{21}
착오에 관한 설명으로 옳지 \uline{않은} 것은? (다툼이 있으면 판례에 의함)
\begin{choices}
  \item 착오를 이유로 의사표시가 취소된 경우 그 취소로 인하여 손해를 입은 자는 불법행위를 이유로 손해배상을 청구할 수 있다.
  \item 금융기관이 아닌 채권자와 체결한 보증계약에서 주채무자의 신용상태나 자력에 관한 보증인의 착오는 중요부분에 관한 것이 아니어서 보증인은 보증계약을 취소할 수 없다.
  \item 매수인이 목적물을 시가보다 고액으로 매수한 경우에 원칙적으로 매수인은 시가의 착오를 이유로 취소할 수 없다.
  \item 경계선을 침범하였다는 상대방의 강력한 주장에 의해 착오로 그간의 경계 침범에 대한 보상금 내지 위로금 명목으로 금원을 지급한 경우에는 법률행위의 중요부분의 착오로서 취소할 수 있다.
  \item 甲이 채무자란이 공란으로 된 근저당권설정계약서를 제시받고 그 채무자가 乙인 것으로 알고 근저당권설정자로 서명날인을 하였는데 그 후 채무자가 丙으로 되어 근저당권설정등기가 경료된 경우 법률행위의 중요부분에 대한 착오가 존재한다.
\end{choices}
\end{question}
\answer{1}
